\section{Threats to validity}
\label{sec:threats}

\paragraph{Virtualised platforms, one real slow device.} All measurements except those
of \S\ref{sec:realdisk} ran in WSL2, a Hyper-V virtual machine, on one NVMe SSD; we did
not measure a hard disk or an SD card. The one slow device we measured, network block
storage under an OpenStack virtual machine, differs from the laptop in kernel (7.0),
interpreter (3.14.4) and processor, served other programs during the runs, cached
recently written data below the guest, and was measured with three repetitions and no
global cache drop, so its interpreter and tools were never cold. Process creation, system
calls and page-cache behaviour are those of a real Linux kernel, so the CPU-bound
costs that dominate Phases~1 and~2 should carry over to bare metal, but their
absolute values will differ. ``Cold'' runs drop the guest's caches and remount the
file system, and the loop devices bypass the guest page cache, but the Windows host
may still cache the virtual disk; the cold timings on the plain loop device are
therefore lower bounds on device time. A calibration with \texttt{fio} (4\,KiB random
direct reads over 8\,GiB written just before) gives 0.34\,ms per read at queue depth 1
and 32{,}000 reads per second at depth 16, the same in a second pass over the same
offsets; from inside the guest we cannot tell how many of these reads the host's cache
served. Our conclusions about slow devices rest on
request counts and on injected delays, which the host cache does not affect.

\paragraph{Emulated devices.} dm-delay and the one-request-at-a-time NBD device bound
device concurrency from both sides but model neither seek distance nor command
reordering. A real disk with native command queueing sits between the two bounds.

\paragraph{Synthetic trees and traces.} The controlled sweeps use freshly created
trees with one extent per file unless stated otherwise. Aged file systems scatter
inodes and directory blocks more, which increases cold-cache device reads per file
(our \emph{scattered} layout approximates this). The real traces come from six
workloads on one machine.

\paragraph{Implementation language.} All rungs are Python, as the artifact is; the
cost of a Python-level loop iteration and of \texttt{fork}/\texttt{exec} are part of
what we measure. A C implementation of either design would shift the constants but
not the subprocess and tree-walk counts that dominate upstream.

\paragraph{One code version.} We measured the fork at commit c162012 and upstream at
796df86. Later changes to either could change the results.
